Die Modernisierung veralteter, sicherheitskritischer Elektronik ist oft kostengünstiger und umweltfreundlicher als ein vollständiger Austausch – vorausgesetzt, dass die funktionale Gleichwertigkeit des Retrofits gegenüber dem Original nachweisbar ist und es nach aktuellen, strengeren Normen erneut zertifiziert werden kann. Diese Fallstudie behandelt diese Herausforderung am Beispiel eines Gleitschutzmoduls für den europäischen öffentlichen Nahverkehr.

Die bisherige 8-Bit-Controller-Platine (basierend auf \gls{mcs48}) wurde durch ein energiesparendes Flash-\gls{fpga} (\texttt{Actel A3P1000}) ersetzt, das die ursprüngliche \gls{cpu}- und \gls{io}-Logik emuliert, ergänzt um einen 32-Bit-Arm\textsuperscript{\textregistered}-\gls{mc} (\texttt{STM32F401RET6}) für erweiterte Protokollierung. Ein zentrales Ziel war die Entwicklung eines methodischen Test-Frameworks, das sich für eine breitere Anwendung bei der Nachrüstung sicherheitskritischer Systeme eignet. Kern dieser Methodik ist eine zweistufige Verifikationsstrategie:
\begin{enumerate*}[label=(\arabic*)]
	\item \textit{Komponenten- und Schnittstellentests:} Deterministische \gls{hdl}-Simulation und Hardware-Tests (Boundary-Scan, Timing) überprüfen die logische Gleichwertigkeit und das Schnittstellen-Timing (z.B. Watchdog, Frequenzsweep).
	\item \textit{Systemtests:} Szenariobasierte, nicht-deterministische Testkampagnen an einem speziellen Prüfstand~\cite{dolata2023} validieren integriertes Verhalten (Fehlerspeicher, Selbsttests, Diagnose).
\end{enumerate*}
Die Ergebnisse wurden auf die Normen \gls{iec61508}~Ed.~2, \gls{en50129}:2018, \gls{en50155}:2021, \gls{en50716}:2023 sowie \gls{ieee1012}:2024 abgebildet. Obwohl vollständige Hardware-in-the-Loop (\gls{hil})-Teststände außerhalb des Projektumfangs lagen, entspricht das implementierte Konzept den \gls{hil}-Abdeckungskriterien, wodurch zukünftige Integrationen erleichtert werden.

Lebensdaueruntersuchungen bestätigen, dass die Modernisierung bestehender Elektronikkomponenten in Bahnfahrzeugen die Betriebskosten erheblich reduziert und den CO\textsubscript{2}-Fußabdruck gegenüber Neuprodukten deutlich senkt. Unsere Analyse bestätigt, dass Retrofit-Platinen vergleichbare wirtschaftliche und ökologische Vorteile bieten. Erfolgreiche Implementierungen durch große europäische Bahnbetreiber validieren zudem die allgemeine Anwendbarkeit dieser rigorosen Verifikationsmethodik.